\chapter{Filling area and homotopy classes of loop families}
\label{ch:filling-width-v4}

The extinction argument needs a numerical invariant which survives a surgery
event. A single loop is not enough: the relevant class is represented by a
two-sphere's worth of loops, and the number assigned to a family is the
largest least area of a disk spanning one of its members. This chapter
constructs that invariant on each selected component of the surgery flow.
First we retain the component and its fundamental-group information through
surgery; then we identify the loop-space class; only afterwards do we take
disk-area infima and family and class infima.

Throughout, $M$ denotes a Hausdorff, second-countable smooth three-manifold
without boundary, and $g$ its Riemannian metric. A surgery slice is denoted
$M_t$. A selected component is a compact connected smooth manifold $C$
with a specified smooth open embedding $\iota:C\to M_t$ whose image is a
connected component, and a specified point $x\in C$. Metrics on $C$ are
pullbacks by this inclusion.

\section{Spheres at a surgery event}
\label{sec:filling-surgery-spheres}

Let $S\subset M$ be a smoothly embedded two-sphere. It is null-homotopic when
its parametrization $S^2\to M$ is null-homotopic. The separation statement used by
surgery is:

\begin{proposition}[Null embedded spheres separate]
\label{prop:filling-sphere-separation}
\uses{prop:foundations-second-homology}
If $M$ is compact, connected and smooth of dimension three and $S$ is a
smoothly embedded null-homotopic two-sphere, then $M\setminus S$ is nonempty and is
not preconnected. Hence it has at least two connected components.
\source{MT2007}\dcref{Proposition 15.12 and Remark 15.13, p. 365}
See Chapter~\ref{ch:foundations-v4} for the topological background.
\lean{PoincareMT.repairedSphereSeparation}
\end{proposition}

\begin{proof}
Use integral singular homology. Since the inclusion of $S$ is null-homotopic,
it sends its fundamental class to zero in $H_2(M;\mathbb Z)$. Exactness for
the pair $(M,S)$ therefore supplies $a\in H_3(M,S;\mathbb Z)$ with
$\partial a=[S]$. There is no choice of orientation on $M$ in this step.

For $p\notin S$, restrict $a$ to the local group
$H_3(M,M\setminus\{p\};\mathbb Z)\cong\mathbb Z$. Whether this class is
twice another class is independent of the sign of the isomorphism.
This parity is locally constant on $M\setminus S$: in a coordinate ball
disjoint from $S$, restriction from a compact convex neighborhood to any
one of its interior points is an isomorphism in degree three. A common
restriction of $a$ consequently gives the same parity at all sufficiently
nearby points. If $M\setminus S$ were connected, the parity would be the
same everywhere on it.

We show that it differs on the two sides of a small piece of $S$. A
straightening chart takes that piece to $\mathbb R^2\times\{0\}$.
Choose a closed cylinder $K=\overline D_r^2\times[-r,r]$ inside the chart
and put
\[
 Q=(\mathbb R^2\times\{0\})\cup K^c,\qquad
 p_+=(0,r/2),\quad p_-=(0,-r/2).
\]
Restriction and excision take $a$ to $b\in H_3(\mathbb R^3,Q;\mathbb Z)$.
Restriction to the two punctures identifies this group with
$\mathbb Z\oplus\mathbb Z$. Excision identifies the support groups of
the two convex half-cylinders with those at $p_+$ and $p_-$, respectively.
Under these coordinates the map
\[
 H_3(\mathbb R^3,K^c;\mathbb Z)\longrightarrow
 H_3(\mathbb R^3,Q;\mathbb Z)
\]
has image a primitive signed diagonal: restriction to either point is an
isomorphism, so a generator maps to $(1,1)$ or $(1,-1)$ after choosing
the sign of the first coordinate.

If the two coordinates of $b$ have equal parity, subtract a suitable
multiple of this diagonal generator. Both remaining coordinates are even,
so $b=d+2c$ with $d$ in the displayed image. The boundary map for
$K^c\subset Q\subset\mathbb R^3$ kills $d$. Thus its value on $b$ is
even in $H_2(Q,K^c;\mathbb Z)$. Naturality of this boundary map and
excision identify that value with the restriction of $[S]$ to the local
disk in $S$. This is a generator of the local integral two-dimensional
homology group and is not even, a contradiction. Hence the complement
is disconnected. The same straightening chart contains points of
positive and negative height, proving that the complement is nonempty.
\end{proof}

For the group calculation, a finite surgery reconstruction starts with a
clopen disjoint union of the pieces and then performs finitely many connected
sums. If $P$ is a surviving piece and $A$ is the parent carrier, van Kampen on
the two punctured summands and their spherical collar gives, for each $x\in P$,
a point $y\in A$ and maps
\[
 q_x:\pi_1(A,y)\longrightarrow\pi_1(P,x),\qquad
 i_x:\pi_1(P,x)\longrightarrow\pi_1(A,y),\qquad q_x i_x={\rm id}.
\]
The factor datum retains the kernel of $q_x$ as a subgroup of the
parent group; it does not identify it with the other summand's subgroup. The
initial clopen union contributes a based group isomorphism; composing the
retractions and sections through the finite connected-sum sequence gives the
displayed maps for the complete reconstruction. The basepoint $y$ is part of
the construction and is not silently identified with $x$.

We spell out the two uses of van Kampen, including basepoints inside
the removed ball. A surgery ball has a radius-two coordinate chart,
and the removed closed ball has radius one. Its open complement and
the radius-two chart cover the original summand. Their intersection
is $S^2\times(1,2)$, hence simply connected, and the chart is
contractible. Filling in this ball therefore induces a fundamental-group
isomorphism at any point of the complement. If the original basepoint
lies inside the removed ball, join it inside the radius-two chart to
a point of radius $3/2$. Transport along this path first moves its
group to a point of the complement. This argument does not assume that
other components of the summand are connected to the surgery ball.

For the connected sum, use the open punctured first summand and the
union of the punctured second summand with the full connecting collar.
Their overlap is the negative half-collar $S^2\times(-1,0)$.
The retraction form of van Kampen, obtained by sending the other
factor to the identity, makes the first open-set group a retract of
the whole group. Interchanging the summands uses the positive
half-collar. Compose each retraction with the ball-filling isomorphism
and the chosen basepoint path to obtain $q_x$ and $i_x$ above.

The initial clopen disjoint union gives the group of each piece as
the group of its own region: a loop or homotopy based in that region
stays there by connectedness of its parameter domain. If a reconstruction
step has retraction $r:G\to H$ and section $j:H\to G$, and the
earlier steps have retraction $r':H\to K$ and section $j':K\to H$,
the composite maps are $r'r$ and $jj'$. Their composite is the identity
on $K$. Induction along the finite reconstruction gives the asserted
maps without assuming connectedness of intermediate carriers. Their
surjectivity and injectivity follow from $q_xi_x=\mathrm{id}$.
The kernel recorded is precisely $\ker q_x$.

\begin{proposition}[Factor and injection effects]
\label{prop:filling-group-effects}
\uses{prop:filling-sphere-separation,thm:surgery-local-topology}
For every finite connected-sum surgery conclusion, each surviving piece has
the factor map, named kernel, and injective section just described. Along a
component path on $[0,T]$ with finitely many surgery times, if the selected
initial fundamental group is trivial, then every selected group is trivial.
\source{MT2007}\dcref{Proposition 15.3, pp. 357--358}
\dcref{Definition 18.2 and Lemma 18.3, pp. 419--422}
\lean{PoincareMT.repairedSurgeryGroupEffects}
\end{proposition}

For persistence, let $E_s$ be the finite set of events at or before $s$. Induct
on $|E_s|$. If $E_s$ is empty, the regular interval supplies a group equivalence
from time zero to $s$. Otherwise take the largest event $e\leq s$. Its stored
predecessor time $e^-<e$ has strictly fewer events, so the induction
hypothesis makes the predecessor group trivial. The survivor injection into
the parent then makes the post-event selected group trivial; the regular
equivalence from $e$ to $s$ finishes the step. The predecessor need not be the
immediately preceding event; strict inequality and finiteness are the only
chronology facts used in this group induction.

\begin{proposition}[Surviving children]
\label{prop:filling-child-components}
\uses{prop:filling-group-effects}
Assume the parent group at each survivor basepoint is subsingleton. Every
surviving piece has subsingleton fundamental group and is simply connected;
the surgery certificate also supplies a point and a specification witnessing
that its survivor region is a connected component of the post-surgery carrier.
\source{MT2007}\dcref{Proposition 15.3, pp. 357--358}
\lean{PoincareMT.repairedChildComponents}
\end{proposition}

Indeed the section $i_x$ is injective, so a subsingleton parent group forces a
subsingleton piece group. A connected smooth manifold is locally path
connected; therefore path connectedness together with the subsingleton group
implies simple connectedness. The component point and its specification are
retained as fields of the child datum. This qualification is essential when
the later width family is transported: its carrier, basepoint and metric are
the selected survivor, not an isomorphic manifold chosen afterwards.

\section{Following components through a finite event set}
\label{sec:filling-component-ancestry}

Fix the controlled flow of Chapter~\ref{ch:global-surgery-v4} with compact
simply connected initial manifold. On a bounded time interval $[0,T]$
its surgery times form a finite set. We need a component path ending at
each prescribed point of $M_T$. Such a path selects a component model
$C_t$ at every time, with regular transports equal to restrictions of the
flow's own slab diffeomorphisms. At an event $s$, it selects the actual
surviving piece containing $C_s$ and an interior retained point in the
incoming component $C_{s^-}$. The time $s^-$ is the reference time stored
in the event, not a newly chosen limiting slice.

There is a geometric chronology assertion behind this definition.
There is no surgery time in $(s^-,s)$. Suppose instead that $u$ is such
a time. Choose $b\in(s^-,u)$. The event's pre-surgery flow has a fixed
compact carrier and is smooth through $u$, so its curvature energy is
bounded on $[b,u]$ times that carrier. The actual pre-surgery
identifications are isometries and transfer this bound to the ordinary
flow before $u$. The singular-time maximality condition at $u$ gives
times arbitrarily close to $u$ from below with curvature exceeding any
specified bound. Choose one greater than $b$ and greater than all earlier
surgery times below $u$; finiteness permits the latter choice. This
contradicts the transferred bound. The required gap is therefore a
consequence of the geometric flow conditions.

Next, every outgoing component contains the image of a point in the
interior of the retained incoming region. The surgery reconstruction
first gives a retained point in that component, possibly on its boundary.
In the boundary case the event's boundary correspondence places its image
on a cap. The cap's connected ball model also contains a negative
half-neck. Under the inverse regular-limit identification, this open
half-neck lies in the incoming retained interior. Its image is in the
same outgoing component as the boundary point, providing the required
interior witness. Thus the argument uses the full local cap model, not
only nonemptiness of a closed overlap.

Now induct forward through the finite event set, making the induction
statement for every point in each slice. On an ordinary interval pull
a target point back by the slab diffeomorphism and extend the earlier
component path by that diffeomorphism. At an event choose an interior
retained point $z$ whose image belongs to the target component. Apply
induction at $s^-$ to $z$; between $s^-$ and $s$ follow its pre-flow
worldline, and at $s$ choose the target point. Components, rather than
points, are being followed, so the outgoing point need not equal the
retained image. The gap just proved ensures that the intervening
worldline crosses no event. The inherited earlier event witnesses and
the new retained witness give all the required compatibility conditions.
Empty slices make the pointwise induction statement vacuous.

In the same finite induction, regular diffeomorphisms preserve trivial
fundamental groups and the survivor injections of
Proposition~\ref{prop:filling-group-effects} preserve triviality at each
event. Starting with the initial diffeomorphism, this proves that every
selected component is simply connected. It does not require a later
whole slice to be connected or even nonempty. Since selected component
ranges are open, their terminal ranges cover the compact slice $M_T$;
a finite subcover gives finitely many actual component paths covering it.
The common initial component will be fixed in
Section~\ref{sec:filling-fixed-initial-class}.

\begin{proposition}[Finite component ancestry]
\label{prop:filling-component-ancestry}
\uses{prop:filling-child-components,thm:global-finiteness}
Every point of every nonnegative-time slice of the controlled flow has
a component path from the initial slice, with the regular and event
compatibilities above. Each selected component is simply connected.
Every target slice has a finite cover by terminal ranges of such paths.
\source{MT2007}\dcref{Definition 18.2 and Lemma 18.3, pp. 419--422}
\lean{PoincareMT.repairedFiniteAncestry}
\end{proposition}

\section{Comparison maps on the selected components}
\label{sec:filling-ancestry-transport}

Fix one of these paths and an event $s$. Write $A=C_{s^-}$ and $B=C_s$,
with inclusions $\iota_-:A\to M_{s^-}$ and $\iota_+:B\to M_s$.
If $R$ is the retained incoming set and $r$ its retention map, use the
entire set
\[
 U=\{a\in A:\ \iota_-(a)\in\operatorname{int}R,
                   \ r(\iota_-(a))\in\iota_+(B)\}.
\]
It is open because retention is continuous on $\operatorname{int}R$
and the child component is open. The ancestry witness proves it is
nonempty. Restrict the pre-flow metrics to $A$ and the outgoing metric
to $B$ by their inclusions. These are the metrics used in the surgery
comparison theorem; no independent metric on a diffeomorphic carrier
is introduced.

Every actual central neck sphere that meets $A$ lies wholly in $A$:
its continuous image is connected, and a connected set meeting a
connected component is contained in it. The inverse regular-limit
identification, followed by the inverse component inclusion, takes its
parametrization to a smooth embedded sphere in $A$. Both inverses are
used only on their open domains. The closed simply connected
three-manifold topology result, Corollary~\ref{cor:foundations-homotopy-groups},
gives $\pi_2(A)=0$, so this embedded
sphere is null-homotopic. Proposition~\ref{prop:filling-sphere-separation}
proves that its actual complement in $A$ is disconnected. Finally
$H_3(A;\mathbb Z)\cong\mathbb Z$ supplies the integral orientation
required by the comparison construction.

These facts supply the component hypotheses of the comparison theorem
in Theorem~\ref{thm:comparison-transport}: the specified surviving child, the
nonempty open retained region, the separating neck spheres, the
restricted metrics, and the integral orientation. At the theorem's
stated smallness bounds on the surgery parameters, take its actual
comparison map $f:A\to B$ and its conclusion that the induced map on
$\pi_3$ is bijective. This is a use of that chapter's comparison theorem;
the bijectivity is not inferred merely from a nonempty retained overlap.

For a chosen incoming point $a_0$, the comparison may give a child
basepoint $f(a_0)$ different from the selected point $b_0$. A path $p$
in the connected smooth manifold $B$ joins them. The transported class
is
\[
 \xi\longmapsto p_*f_*\xi\in\pi_3(B,b_0).
\]
Postcomposition sends the constant representative to the constant
representative. Since $f_*$ is bijective, $f_*\xi=0$ implies $\xi=0$.
Since reverse-path transport is inverse to $p_*$ and preserves zero,
$p_*f_*\xi=0$ has the same implication. Thus a nonzero class survives
the event. On a regular slab use the restricted slab diffeomorphism;
its inverse proves bijectivity by composition on representatives, and
the same path argument handles the selected basepoint.

\begin{proposition}[Transport along an actual ancestry path]
\label{prop:filling-ancestry-transport}
\uses{prop:filling-component-ancestry,thm:comparison-transport}
With the surgery comparison hypotheses and the path-induced basepoint
maps, nonzero third homotopy classes persist through the regular and
event maps just constructed on every component path.
\source{MT2007}\dcref{Proposition 15.12 and Remark 15.13, p. 365}
\dcref{Claim 18.16 and Proposition 18.18, pp. 430--432}
\lean{PoincareMT.repairedAncestryTransport}
\end{proposition}

\section{The loop-space class}
\label{sec:filling-loop-class}

Let $\Lambda^1M$ be the space of $C^1$ free loops. In the model used here a
loop carries an extension to a neighborhood of the circle, but the topology
observes only its values and its tangential first derivative on the circle.
Thus convergence means uniform convergence of these two maps into $M$ and
$TM$, respectively. Off-circle values and normal derivatives do not
contribute to this topology.

Fix $I=[0,1]$ and a continuous surjection $q:I^2\to S^2$ whose fibers are
singletons in the interior and the entire boundary over a pole $c_*$. A
raw null family is a continuous map $F:S^2\to\Lambda^1M$ such that each
$F(c)$ is null-homotopic in $M$. This condition is pointwise and does not
choose contractions continuously in $c$. A normalized regular family at
$x$ has $F(c_*)=\gamma_x$, the constant loop at $x$, so $F\circ q$
represents a based second homotopy class. It also carries extensions
continuous jointly in the sphere and loop parameters, with jointly
continuous first derivatives on the circle. The fixed quotient retains
the orientation of the parameter sphere when assigning this class.

For every compact connected smooth three-manifold with $\pi_2(M,x)=0$
there is a group isomorphism
\[
 \Phi_{M,x}:\pi_2(\Lambda^1 M,\gamma_x)\xrightarrow{\cong}\pi_3(M,x).
\]
Moreover: (i) a loop lies in the component of $\gamma_x$ exactly when it
is null-homotopic; (ii) every class has a normalized regular family using $q$;
(iii) two normalized regular families are freely homotopic exactly when their
based classes agree, with a homotopy through regular records; and (iv) the
absolute and relative square models agree, where the relative condition
fixes the edges $u=0$, $v=0$, and $v=1$ at $\gamma_x$, and permits only
constant loops on $u=1$.

Here are the distinct comparisons needed for this assertion. First, the
value map $\Lambda^1M\to C(S^1,M)$ induces isomorphisms on positive
cubical homotopy groups. Given a compact-parameter continuous family of
circle maps, finite-chart approximation makes it $C^1$ in the circle
variable while fixing the constant boundary family. To make this
approximation, cover the compact parameter domain by finitely many
bump plateaus mapping into target charts. Replace the map successively
in these charts, preserving its already smooth locus and the full
inverse image of the prescribed constant value. Before each step,
compactness of the remaining bump supports gives a positive tolerance
preserving every remaining chart condition. Use less than half the
remaining error at that step. The finite composition is smooth on
the covering plateaus, fixes the required constant parameters, and
has the prescribed uniform value error. A uniform
neighborhood of the diagonal admits a local contraction; it joins the
approximation to the original family, still fixing the boundary.
This proves surjectivity. For injectivity, approximate the whole
homotopy cylinder in one additional cube dimension, fixing its
constant boundary parameters. Its endpoint approximants need not be
the original $C^1$ families, but are uniformly close to them. The
same contraction joins the first original family to its approximant
and the last approximant to the last original family through $C^1$
loops. Concatenating these two homotopies with the approximating
cylinder gives a based homotopy with the exact original endpoints.
A path to a constant loop evaluates to a null
homotopy, and the same relative approximation converts a continuous null
homotopy into a path of $C^1$ loops. Connectedness moves its final constant
point to $x$, proving (i).

The regular-family construction does not require differentiating in
the sphere parameter. Replace each annular loop extension by its
radial extension $\widehat\gamma(z)=\gamma(z/|z|)$. On $|z|=1$ this
has the same values and tangential first derivative as $\gamma$.
Its derivative in any plane direction $v$ is
\[
 d\widehat\gamma_z(v)=
 \langle T_z,v\rangle\,d\gamma_z(T_z),
\]
where $T_z$ is the oriented unit circle tangent. Thus a continuous
family in the intrinsic first-jet topology has a jointly continuous
radial extension and jointly continuous ambient first derivatives
on the circle. Radialization fixes the whole constant-loop
extension. Retaining the original extension at homotopy time zero
and using its radial extension at every positive time is continuous
in this topology, because all observed values and tangential jets
are unchanged.

A based square descends through $q$ to a pole-normalized sphere
family. Each of its loops is connected to the pole loop by a path
in the parameter sphere, and therefore is null. Radialization gives
the asserted regular representative without changing its based
class. For an arbitrary raw null family, choose a path in the loop
space from its pole value to $\gamma_x$. Transport the boundary of
its square along that path; the accompanying homotopy descends
through $q$ to a free sphere homotopy. Its final family is normalized
at $x$, and can now be radialized. Finally, a homotopy between
normalized based classes descends to a pole-fixed sphere homotopy.
Radialize its interior time slices while retaining the two given
endpoint extensions. Equality of the observed first jets proves
continuity even at those endpoints. This establishes raw
regularization and regular homotopies with the original endpoints.

Second, inclusion of based loops in $C(S^1,M)$ induces an isomorphism on
$\pi_2$ when $\pi_2(M,x)=0$. One can see this directly by transposing
parameters: a based two-cube of free loops is a free loop in the space
of based two-cubes in $M$. This latter space is path connected by
$\pi_2(M,x)=0$. A path from the value at the circle basepoint to the
constant cube turns the free loop into a based loop, proving
surjectivity. Two based loops freely homotopic in this cube space have
classes differing by conjugation. Its fundamental group is abelian,
by interchange of the positive-dimensional cube coordinates, so those
classes agree. This proves injectivity. Finally interchange of the
circle and square coordinates gives the cubical adjunction
$\pi_2(\Omega M)\cong\pi_3(M)$. Composing these comparisons defines
$\Phi_{M,x}$.

Free homotopy of sphere families needs a further argument. Let
$\widetilde M\to M$ be the path-class universal cover. Covering
homotopy lifting identifies higher based homotopy groups, so
$\pi_2(\widetilde M)=0$. If $\widetilde M$ is noncompact, its ordinary
integral third homology vanishes. Indeed a finite singular cycle has
compact support and hence misses some point of $\widetilde M$. Its
local homology coefficient is zero there. Simple connectedness supplies
coherent integral local orientations; the coefficients of this cycle
are locally constant, hence constant, on the connected manifold. All
are therefore zero, and the compact-support detection theorem forces
the original homology class to vanish. Hurewicz in degree three now
gives $\pi_3(\widetilde M)=0$, and consequently $\pi_3(M)=0$.

If the cover is compact, use the finite triangulation of the base from
Proposition~\ref{prop:foundations-triangulation} and
lift each characteristic simplex to the cover. The normalized integral
chains on these lifts form a finite free complex; characteristic
projection to singular chains induces homology isomorphisms and commutes
with deck maps. A nonidentity deck map fixes no point: a fixed point
would make it the identity by uniqueness of lifts on a connected cover.
It therefore fixes no lifted simplex, since evaluating a fixed lift
would give a fixed point. Its chain matrix has zero diagonal in every
degree. The alternating chain trace is zero. The trace identity for a
finite free complex equates this to the alternating homology trace.
The simply connected closed three-manifold calculation gives integral
homology $\mathbb Z$ in degrees zero and three and zero otherwise.
Thus
\[
 0=1-\operatorname{tr}(d_*|H_3(\widetilde M;\mathbb Z)),
\]
and the action on $H_3\cong\mathbb Z$ is the identity. The identity
deck map has the same conclusion. Naturality of Hurewicz and of
basepoint transport implies that each deck map on $\pi_3$ agrees with
transport along a path from a point to its deck image.

To apply this, pull a free homotopy of normalized sphere families back
through $q$ and take loop values. Lift the resulting parameter homotopy
to $\widetilde M$. Its two normalized endpoints differ by a deck map
and the path followed by the basepoint. The preceding calculation makes
these changes agree on $\pi_3$. Descend by the covering and use
injectivity of the $C^1$ value comparison to obtain equality of the
original based family classes. Conversely, a based cube homotopy
descends through $q$ to a free sphere homotopy. The same relative
construction retains the given regular endpoint families. This explains
why the free-class assertion uses closed three-manifold topology, in
addition to the based adjunction.

For completeness, the lifted characteristic comparison in the compact
argument can be constructed on the given triangulation; an arbitrary
isomorphism of homology groups would not preserve the deck action.
Write the finite triangulation as the order complex of a finite poset
$J$. For a selected vertex set $s\subset J$, let $D_s$ be the lifted
simplicial subcomplex whose vertices lie in $s$, and let
\[
 W_s=\left\{z\in |J|:\sum_{i\in s}z_i>0\right\}.
\]
Here the $z_i$ are nonnegative barycentric coordinates of sum one,
with support a chain. The characteristic maps of $D_s$ land in
$p^{-1}(W_s)$ and define a natural comparison to its singular chains.

Induct on the number of selected vertices. The empty case has zero
chains and empty target. Choose a minimal vertex $v\in s$ and set
\[
 a=s\setminus\{v\},\qquad b=\{j\in s:v\leq j\},\qquad
 l=\{j\in s:v<j\}.
\]
Then $D_a\cup D_b=D_s$ and $D_a\cap D_b=D_l$.
Also $W_a\cup W_b=W_s$ and $W_a\cap W_b=W_l$.
For the latter intersection, a point with positive $v$-coordinate
and another positive selected coordinate has that second coordinate
strictly above $v$: the support is a chain and minimality excludes a
distinct selected vertex below $v$. If its positive $b$-coordinate
is already above $v$, it directly belongs to $W_l$.

The comparison for $b$ is a cone calculation. First normalize the
selected barycentric weights, deforming $W_b$ into the supported
subcomplex, and then contract to $v$, which lies below every vertex
of $b$. This is a strong contraction fixing $v$. Lift it through the
covering. The covering over $W_b$ is a product with the discrete
fiber over $v$, and the lifted contraction retracts it to that fiber.
On simplicial chains, adjoining the least vertex gives the cone
contraction, one copy for each point of the same fiber. The actual
characteristic projection commutes with these augmentations, so it
induces the same degree-zero group and zero positive homology.

The sets $a$ and $l$ have fewer vertices, so induction handles them.
The simplicial union sequence and singular Mayer--Vietoris sequence
for the open cover $p^{-1}(W_a),p^{-1}(W_b)$ commute with the actual
characteristic maps. Exactness and the five lemma give the comparison
for $s$. With all vertices selected, $W_J=|J|$ and $D_J$ is the
entire lifted model. This proves the comparison used in the deck trace
argument, with its naturality retained. It is the finite-order-complex
form of the simplicial/singular comparison in
\cite[Theorem 2.27, pp.~128--130]{Hatcher}.

The system is coherent under smooth postcomposition. If $f:M\to N$ is
smooth, its loop map postcomposes the chosen loop extensions, not just loop
values. It sends constants to constants and retains $q$. When both
manifolds satisfy the hypotheses above, it satisfies
\[
 \Phi_{N,f(x)}(f_*\alpha)=f_*\Phi_{M,x}(\alpha).
\]
This naturality lets one class label cross a comparison map at a surgery
event.

Basepoints in the same connected smooth component can be joined by paths.
For a path $p:x\to y$, whiskering a cube representative by $p$
defines $p_*:\pi_k(M,x)\to\pi_k(M,y)$. The construction is on representatives
and descends because homotopic representatives remain homotopic after
whiskering. Reverse-path and concatenation homotopies prove
\[
 (p\mathbin{\cdot}q)_*=q_*p_*,\qquad (p^{-1})_*p_*={\rm id},
\]
and the maps preserve identity and multiplication for $k>0$. Continuous maps
commute with this transport. Thus a selected child point can be used without
an unrecorded algebraic isomorphism.

\begin{proposition}[Coherent loop classes]
\label{prop:filling-loop-identification}
\uses{cor:foundations-homotopy-groups,prop:foundations-cw-equivalence}
There is one quotient and one identification system with these properties. On
a selected simply connected child component, every specified $\pi_3$ class has a
normalized regular representative with the same metric, basepoint and
quotient parameter. For each such compact Riemannian component there is
a positive length threshold such that, if every loop in a family has
length below that threshold, its represented class is the identity.
\source{MT2007}\dcref{Claim 18.16 and Definition 18.17, p. 430}
\dcref{Lemma 18.27, p. 434}
\lean{PoincareMT.m59LoopClassesAndComponentTopology}
\end{proposition}

The short-loop clause is not a statement about arbitrary short loops in a
noncompact manifold: compactness, connectedness and $\pi_2(M,x)=0$ are
hypotheses of the raw family theorem.

\subsection{The relative square}
\label{sec:filling-relative-square}

The relative condition stated above is useful because it allows one edge
to move through constant loops. We justify both directions of its
comparison with the fully based square. Let $F:I^2\to\Lambda^1M$
satisfy that condition and fix a circle point $z_0$. On the boundary of
the square define
\[
 K(s;u,v)=\gamma_{F((1-s)u,v)(z_0)},\qquad 0\leq s\leq1.
\]
At $s=0$, this agrees with $F$ on the boundary, since every boundary
loop of $F$ is constant. It fixes the three prescribed edges at
$\gamma_x$, and at $s=1$ it is $\gamma_x$ on all four edges.
The homotopy-extension property for the boundary of a square extends
$K$ to a homotopy of $F$ through relative squares. Its final square is
fully based. This surjectivity argument needs no vanishing hypothesis.

For injectivity, suppose a relative homotopy $H(t;u,v)$ joins two fully
based squares. Evaluating its moving edge gives
\[
 B(t,v)=H(t;1,v)(z_0):I^2\longrightarrow M.
\]
This is $x$ on the entire boundary: the endpoint squares are based,
and the bottom and top edges stay fixed. If $\pi_2(M,x)=0$, choose a
based null homotopy $J(s;t,v)$ of $B$, from $B$ at $s=0$ to $x$ at
$s=1$. On the side boundary of the homotopy cube use the deformation
\[
 (s;t,u,v)\longmapsto\gamma_{J(1-u(1-s);t,v)}.
\]
At $s=0$ it is the old side trace: on $u=1$ it is the constant loop at
$B(t,v)$, and on the other three edges it is $\gamma_x$. It leaves both
endpoint squares fixed because $J$ is based in the $t,v$ square. At
$s=1$ the side trace is constant. Homotopy extension relative to the
two endpoint squares therefore produces a homotopy fixing all four
edges. Hence two absolute classes agree exactly when their relative
classes do. The construction also commutes with postcomposition,
because evaluation and constant-loop inclusion do.

\section{Short loops and actual filling area}
\label{sec:filling-short-loops}

For a $C^1$ loop $\gamma$ let $\ell_g(\gamma)$ be the integral of metric speed over the
fixed loop circle. We first construct one contraction of nearby pairs,
so its application to a family makes no discontinuous choice of chart.
Take finitely many smooth bumps $b_i$ supported inside coordinate charts
$\varphi_i$, with their plateau neighborhoods $b_i=1$ covering $M$.
The $i$th operation, wherever both points lie in its chart, is
\[
 C_i(t,p,q)=\varphi_i^{-1}
 \bigl((1-tb_i(p))\varphi_i(q)+tb_i(p)\varphi_i(p)\bigr);
\]
outside that chart pair it is the second-point projection. Near a
diagonal point in the support, the segment remains inside the open
chart target, so this formula is smooth there. Near a point outside
the support the weight is zero and the operation is the projection.
Thus every $C_i$ is smooth at every diagonal point. Compose the
operations in a fixed order, keeping $t,p$ fixed and changing $q$.
Each starts at $q$ and fixes the diagonal at every time. For each $p$,
a plateau operation sends a sufficiently nearby pair to $p$ at time
one, and every subsequent operation fixes $p$.

For any prescribed finite differentiability order $k$, the set where
this finite composition is $C^k$ is open and contains
$[0,1]\times\operatorname{diag}M$. The tube lemma and compactness
give one open diagonal neighborhood $U$ on which it is $C^k$ for
every $t\in[0,1]$. Shrink $U$ so the time-one identity also holds
throughout it. Only a fixed finite order, chosen for each use, is
claimed on this neighborhood.

A compact Riemannian manifold has a positive uniform radius $\rho$
whose close pairs belong to $U$. Since distance from a loop's initial
point is bounded by its length, every loop shorter than $\rho$ gives
pairs in $U$. The resulting contraction is based at its initial value
and has a continuous time parameter;
applying it to a raw family $F$ gives a homotopy to the family of constant
loops at $F(c)(z_0)$, where $z_0$ is the fixed circle basepoint. The
endpoint identities are pointwise in the loop circle, so they define a
homotopy in $\Lambda^1M$. This terminal family may vary with $c$.
For a connected $M$ with $\pi_2(M,x)=0$, its evaluation map
$c\mapsto F(c)(z_0)$ is null-homotopic and can be contracted to $x$.
Applying the constant-loop inclusion to that homotopy contracts the
terminal family to the constant family at $x$. For a normalized based
family, the contraction fixes its prescribed constant-loop boundary,
so the same argument gives the zero based class.

Write the contraction as $C(t,p,q)$ with $C(0,p,q)=q$,
$C(1,p,q)=p$, and $C(t,p,p)=p$. The disk estimate uses
$C(\theta(r),\gamma(z_0),\gamma(z))$, where the smooth profile $\theta$
is one for $r\leq1/2$ and zero on the boundary $r=1$.
The value at the center is $\gamma(z_0)$. Uniform bounds on radial
and angular derivatives follow from the finite coordinate cover. In polar
coordinates the area density satisfies
\[
 J_F(r,\theta)\leq
 \|\partial_rF(r,\theta)\|_g\,
 \|\partial_\theta F(r,\theta)\|_g.
\]
More precisely, let $A$ bound the time derivative of the contraction,
$B$ bound its derivative in the moving point, and $H$ bound the derivative
of the radial time profile. These bounds are uniform on the compact
sets selected from the finite coordinate cover. Polar integration then
gives
\[
 \operatorname{Area}_g(D_\gamma)\leq K\ell_g(\gamma),\qquad K=HAB\geq0.
\]
The angular speed integrates to the loop length, and the radial integral
is bounded by the fixed profile constant. The center causes no
singularity because the profile is constant there. For a given
$\eta>0$, choose
\[
 \zeta=\min\left\{\rho,\frac{\eta}{4},\frac{\eta}{K+1}\right\}.
\]
Then $0<\zeta<\eta/2$, and $\ell_g(\gamma)<\zeta$ implies
$K\ell_g(\gamma)<\eta$. The resulting disk is $C^1$, hence Lipschitz
on its compact domain, and has exactly the prescribed boundary
parametrization. This establishes a uniform linear area estimate.

For a null loop, an arbitrary continuous null extension is first clamped to
be regular near the boundary. Relative smoothing on a compact inner disk
produces a global $C^1$ extension without changing boundary values. Its metric
Jacobian is integrable and nonnegative, so it is admissible. Define
\[
 A_g(\gamma)=\inf\{ {\rm Area}_g(D):D\text{ is a Lipschitz spanning disk for }\gamma\}.
\]
Here an admissible disk has boundary $D(z)=\gamma(\rho(z))$ for a
specified circle homeomorphism $\rho$, which need only be continuous.
Its area is the integral of its metric Jacobian, evaluated on its
almost-everywhere differentiability locus. The short disks constructed
above use $\rho=\mathrm{id}$. The full infimum initially allows all
these homeomorphisms.
The set is nonempty for null loops and bounded below by zero. Hence
$A_g(\gamma)\geq0$ and $A_g(\gamma)\leq\operatorname{Area}_g(D)$ for every
admissible $D$. The infimum need not be attained; for every $\epsilon>0$
there is an admissible $D$ with
$\operatorname{Area}_g(D)<A_g(\gamma)+\epsilon$.

Filling area is invariant under boundary homeomorphisms relating two
$C^1$ loops. Keep the disk map and compose its boundary label with the
given homeomorphism; the inverse relabeling proves equality of the
entire area ranges. This does not yet give exactly parametrized disks.

\begin{lemma}[Exact boundary parametrization without area cost]
\label{lem:filling-boundary-regularization}
Every admissible disk for $\gamma$ has a Lipschitz replacement of the
same area whose boundary is exactly $\gamma$.
\lean{PoincareMT.m60Disk_regularize_boundary}
\end{lemma}

\begin{proof}
Lift the boundary homeomorphism to a homeomorphism $H:\mathbb R\to
\mathbb R$ in the angular variable, of period $2\pi$ up to sign.
If it reverses orientation, first reflect the parameter disk. Reflection
preserves its area and changes the sign, so assume
$H(t+2\pi)=H(t)+2\pi$ with $H$ increasing.

Let $f(t)=\gamma(e^{it})$, let $\ell(t)$ be its signed accumulated
length from $0$ to $t$, and let $\beta$ be its natural length
parametrization. Thus $f=\beta\circ\ell$, and $\beta$ is
one-Lipschitz on the interval $\ell(\mathbb R)$. Although $H$ need
not be Lipschitz, both $\ell$ and $a=\ell\circ H$ are Lipschitz:
the former follows from the $C^1$ speed bound of $\gamma$, and the
latter from the Lipschitz bound of the actual disk trace
$f(H(t))=D(e^{it})$. Monotonicity of $H$ identifies the variation of
this trace on an interval with the corresponding increment of $a$.
Both functions increase by the total loop length $L$ when $t$ increases
by $2\pi$, and $\beta(s+L)=\beta(s)$ on its range.

On $1/2\leq r\leq1$ set $u=2r-1$ and define
\[
 A(r,e^{it})=\beta\bigl((1-u)a(t)+u\ell(t)\bigr).
\]
Convexity of $\ell(\mathbb R)$ makes the argument of $\beta$ valid.
The period identities make $A$ single-valued on the annulus. Its inner
trace is $D(e^{it})$ and its outer trace is $\gamma(e^{it})$.
In each angular chart it is controlled by a Lipschitz scalar function
$s$, with $d(A(z),A(w))\leq|s(z)-s(w)|$; compact annular charts give
a finite Lipschitz bound.

This scalar bound also proves zero area without differentiating
$\beta$. At points where $s$ and $A$ are differentiable, every vector
in $\ker ds$ lies in $\ker dA$, by taking directional increments in
the distance inequality. Since the domain is two-dimensional,
$dA$ has rank at most one. Rademacher's theorem gives this conclusion
almost everywhere; a countable angular cover and the measure-zero
boundary circles give zero integral area on the whole annulus.

Copy $D$ into the disk of radius $1/2$ by $z\mapsto D(2z)$ and glue
it to $A$. The inner traces agree. A segment crossing the common
circle proves a global Lipschitz bound from the two bounds on either
side. The two-dimensional Jacobian scaling cancels the change in
parameter area, so the inner disk still has area $\operatorname{Area}(D)$;
the annulus contributes zero. The outer trace is exactly $\gamma$.
\end{proof}

Thus near-minimizers can always be chosen with exact parametrization.
To prove continuity on null loops, fix $\gamma_0$. Use one constant-fixing local
contraction $C$ and a bound $B$ for both of its endpoint derivatives.
Its time derivative vanishes when the endpoints agree. Continuity and
compactness of $[0,1]\times\operatorname{diag}M$ therefore make that
time derivative arbitrarily small on a sufficiently narrow neighborhood
of the diagonal, with the same $C$ and $B$. A small $C^1$ neighborhood
of $\gamma_0$ gives pointwise close loop values and a common bound for
their lengths. The collar
$C(t,\gamma_0(z),\gamma(z))$ consequently has small area: its angular
derivative is bounded by $B$ times the sum of the two loop speeds,
and its time derivative can be made arbitrarily small. Integrating
their product bounds the area by this small constant times
$B(\ell_g(\gamma_0)+\ell_g(\gamma))$.

Attach this parametrized collar to a near-minimizing disk for
$\gamma_0$, shrinking that disk radially and leaving its parametrized
boundary exact, for an upper bound on $A_g(\gamma)$. Start with a
near-minimizer for $\gamma$ and attach the reverse collar for the lower
bound. With each near-minimizer error and each collar area less than
$\epsilon/2$, the two bounds give
$|A_g(\gamma)-A_g(\gamma_0)|<\epsilon$. The proof uses exact boundary maps; a free
reparameterization is not silently replaced by equality of loop values.

The same argument gives the uniform form needed for approximation of
an entire family.
\begin{proposition}[Uniform small collars]
\label{prop:filling-uniform-collars}
Fix a compact smooth Riemannian three-manifold, a finite bound
$\Lambda\geq0$, and $\eta>0$. There is $d>0$ such that if two
$C^1$ loops $\gamma_0,\gamma_1$ satisfy
\[
 d_g(\gamma_0(z),\gamma_1(z))<d\quad\hbox{for all }z,\qquad
 \ell_g(\gamma_0)+\ell_g(\gamma_1)\leq\Lambda,
\]
there is a parametrized Lipschitz collar between them of area less
than $\eta$. Attaching it, in either direction, to any admissible
disk after exact boundary regularization produces an admissible disk
with area less than the old area plus $\eta$. If the loops are null,
$|A_g(\gamma_0)-A_g(\gamma_1)|<\eta$.
\lean{PoincareMT.m60_exists_controlled_local_contraction}
\end{proposition}

Choose the fixed endpoint derivative bound $B$ above and require the
time derivative to be less than
$\eta/[2(B+1)(\Lambda+1)]$. Compactness gives a single $d$ for this
requirement. The derivative-product integral bounds the collar area
by at most $\eta/2$. Exact regularization and disk shrinking preserve
the old area. Taking near-minimizers with error less than $\eta/2$
in both directions proves the strict filling-area difference.
This includes $\Lambda=0$; no positive lower bound for loop length
is required.

\begin{proposition}[Short-loop filling estimate]
\label{prop:filling-short-area}
\uses{prop:filling-uniform-collars,lem:filling-boundary-regularization}
For every compact smooth Riemannian three-manifold and every $\eta>0$ there is
$0<\zeta<\eta/2$ such that
\[
 \ell_g(\gamma)<\zeta\Longrightarrow
 \gamma\text{ is null-homotopic and }A_g(\gamma)<\eta.
\]
For a continuous raw family whose every member is shorter than $\zeta$,
its family maximum is less than $\eta$.
\source{MT2007}\dcref{Lemma 18.27 and corrected Corollary 18.28, pp. 434--435}
\lean{PoincareMT.LoopSpace.small_loop_filling}
\end{proposition}

The family assertion follows by evaluating the pointwise strict bound at
a maximizing parameter, whose existence is proved below. Neither
$\pi_2$- nor $\pi_3$-vanishing is needed to construct these disks.

\section{Area of sphere maps and the two width infima}
\label{sec:filling-width-definitions}

For a sphere map $f:S^2\to M$ use the fixed stereographic parameter
$\sigma:\mathbb R^2\to S^2\setminus\{p_\infty\}$. For a plane map $h$,
with $dh(e_i)$ its two coordinate derivatives, put
\[
 G_h(z)_{ij}=g(dh_z(e_i),dh_z(e_j)),\qquad
 J_h(z)=\sqrt{\max\{0,\det G_h(z)\}},
\]
\[
 {\rm Area}_g(f)=\int_{\mathbb R^2}J_{f\circ\sigma}(z)\,dz,\qquad
 E_g(f)=\frac12\int_{\mathbb R^2}{\rm tr}G_{f\circ\sigma}(z)\,dz.
\]
The Jacobian counts parametrized multiplicity. Cauchy--Schwarz gives
$\operatorname{Area}_g(f)\leq E_g(f)$, with equality for a weakly conformal map, including a zero
scale at branch points. A branched minimal sphere is a smooth nonconstant
weakly conformal map stationary for $E_g$, with finite differential-zero set and
injective differential off that set.

We also need to record the least-sphere conclusion carried by the area
theory. It holds for a compact smooth Riemannian target $X$ of arbitrary
dimension, provided $\pi_2(X,x)\ne0$ at some point. Set
\[
 m=\inf\{\operatorname{Area}(h):h:S^2\to X\text{ is }C^1
                         \text{ and not null-homotopic}\}.
\]
The infimum is over all non-null spheres, not over a prescribed free
homotopy class. The conclusion is that $m>0$ and a smooth branched
minimal sphere attains it.

First identify this area infimum with the infimum of the half-energy
$E$ over smooth non-null spheres. The inequality area at most energy
gives one direction. For the other, approximate the continuous
semidefinite pullback metric $h^*g$ from above by smooth positive
metrics, adding a small multiple of the round metric to control
degeneracy. Uniformize each such metric on $S^2$ and precompose $h$
with the resulting diffeomorphism. Here the classical uniformization
input is that the complex structure induced by any smooth oriented
metric on $S^2$ is conformally equivalent to the Riemann sphere;
compactness excludes the plane and disk alternatives in
\cite[Theorem 1.11.1, p.~14]{MilicicRiemannSurfaces}.
The energy measured in these
coordinates approaches the original area as the positive error tends
to zero. Precomposition by a diffeomorphism preserves non-nullness.
More explicitly, if $q_\epsilon$ is a smooth metric with
$h^*g\leq q_\epsilon\leq h^*g+\epsilon g_0$, and
$\phi_\epsilon:(S^2,g_0)\to(S^2,q_\epsilon)$ is conformal, then
\[
 E_g(h\circ\phi_\epsilon)
 \leq E_{q_\epsilon}(\phi_\epsilon)
 =\operatorname{Area}_{q_\epsilon}(\phi_\epsilon)
 =\operatorname{Area}_{q_\epsilon}(\operatorname{id}).
\]
The last expression tends to $\operatorname{Area}_g(h)$: in the
stereographic chart, polarization gives convergence of every Gram
coefficient, and the determinant densities are dominated by the
integrable energy density of $h$ plus a fixed multiple of the round
density. This also covers the rank-deficient points of $h$.
Finite supported smoothing in target charts then gives a smooth
homotopic map with arbitrarily small extra energy. The smoothing steps
preserve the already smooth regions and divide the value and energy
error budgets among finitely many charts. Thus smooth non-null
competitors have energy arbitrarily close to $m$.

For $\alpha>1$ use the perturbed functional on the unit round sphere
\[
 E_\alpha(h)=\int_{S^2}(1+|dh|^2)^\alpha\,d\mu_0,
 \qquad E_1(h)=4\pi+2E(h).
\]
Here is the variational construction needed for this particular
infimum. Fix $1<\alpha\leq2$ and take a smooth non-null minimizing
sequence $f_j$ for $E_\alpha$. One fixed smooth non-null competitor
gives a finite upper bound, even uniformly for $1\leq\alpha\leq2$,
after discarding initial terms. A smooth embedding of the compact
target in Euclidean space turns the intrinsic energy bound into a
$W^{1,2\alpha}$ bound on its observed coordinate functions. The
two-dimensional Sobolev inequality gives equicontinuity because
$2\alpha>2$. Arzela--Ascoli gives a uniformly convergent subsequence
with continuous target-valued limit $f_\infty$. Nearby maps are
homotopic by the local contraction of pairs in $X$, so this limit is
non-null. No single free class had to be selected before taking the
sequence.

The derivative convergence requires a variational argument. On a
smaller common source and target chart, write the coordinate maps as
$u_j$ and choose a smooth cutoff $0\leq\chi\leq1/2$. Replace the pair
$u_j,u_k$ by
\[
 (1-\chi)u_j+\chi u_k,\qquad
 (1-\chi)u_k+\chi u_j.
\]
The replacements lie in a convex target-chart ball and agree with the
original maps outside the patch. Linear interpolation there makes
each replacement homotopic to its original sphere, hence non-null.
Uniform convexity of the $2\alpha$-power derivative term, and uniform
continuity of the metric coefficients on that ball, give
\[
 c\int\chi\,|Du_j-Du_k|^{2\alpha}
 \leq Q(\delta)\bigl(E_\alpha(f_j)+E_\alpha(f_k)\bigr)
       -2m_\alpha+e(\delta),
\]
where $c>0$, $m_\alpha$ is the smooth non-null energy infimum,
$\delta$ bounds $\|u_j-u_k\|_\infty$ and the resulting coefficient
oscillation, $Q(\delta)\to1$, and $e(\delta)\to0$.
The cutoff-derivative terms have the factor $u_j-u_k$ and enter
$e(\delta)$. First make $\delta$ small and then use convergence of
the minimizing energies. On a disk where $\chi$ is a positive
constant the derivatives are Cauchy in $L^{2\alpha}$. Their limit is
the weak derivative of $u_\infty$, by integration against compactly
supported test functions. This proves strong local convergence of
both the values and derivatives.

Supported variations of $u_\infty$ now inherit local minimality from
the original sequence. Differentiating their density integrals gives
the weak Euler equation. In a conformal source chart with round
factor $\lambda$, it has the covariant form
\[
 \sum_{i=1}^2\nabla_i(w\,\partial_i u)=0,\qquad
 w=\left(1+\frac{\sum_i g_u(\partial_i u,\partial_i u)}{\lambda}\right)^{\alpha-1}.
\]
The same map gains smoothness for $\alpha$ in a fixed sufficiently
small interval above one. To indicate the regularity mechanism,
testing the weak equation with supported difference quotients first
bounds the difference quotients of its weak first derivatives in
$L^2$ on a smaller disk. This produces a weak Hessian; the
two-dimensional Sobolev inequalities then give every finite local
integrability exponent for the first derivatives. Normalize the
target metric at the center by a fixed linear change of coordinates
and shrink the disk. The equation for the weak Hessian is a Laplace
equation with an error bounded by
$16(\alpha-1)|D^2u|$ and a lower term in $L^4$.
Absorption in the interior elliptic estimate gives a
$C^{1,1/4}$ representative of this same map. Differentiating the
equation for its first jet and applying the interior estimate
successively gives every finite differentiability order. The small
interval is chosen before the chart and the map; the chart radii
may shrink afterwards. This is the local regularity argument of
\cite[Proposition 2.3, p.~7]{SacksUhlenbeck1981}, with the metric
normalization making that order of choices explicit.

Once $f_\infty$ is smooth, its weak derivatives equal its classical
derivatives almost everywhere. The strong local convergence gives
$L^1$ convergence of the actual $\alpha$-energy densities. A finite
chart cover therefore proves
$E_\alpha(f_\infty)=m_\alpha$. Thus these are absolute minima among
all smooth non-null spheres, obtained without replacing the original
map during regularity. Choose the interval small enough to lie in
$(1,33/32]$ and choose $\alpha_j\downarrow1$ in it. Denote the
resulting smooth minima by $h_j$.

For every fixed smooth competitor $h$,
$E_{\alpha_j}(h_j)\leq E_{\alpha_j}(h)$ and
$E_{\alpha_j}(h)\to4\pi+2E(h)$. On the other hand
\[
 m\leq\operatorname{Area}(h_j)\leq E(h_j)
       \leq\frac{E_{\alpha_j}(h_j)-4\pi}{2}.
\]
Taking competitors with energy arbitrarily close to $m$ squeezes all
three expressions on the right to $m$. The subtraction of $4\pi$ and
division by two are essential to this comparison.

We next extract a harmonic sphere, retaining enough information for
the area comparison. There are two alternatives. If
$\sup_{j,p}|dh_j(p)|^2<\infty$, the fixed exponent-two energy is
bounded, so compactness gives a uniform limit on the entire sphere.
Its non-nullness is retained by the same nearby-pair homotopy as
above. The local weighted equations give first-derivative compactness
and a smooth harmonic limit; convergence of the energy densities
gives its energy at most $m$.

If the gradients are unbounded, choose a subsequence and actual
maximum points $p_j$ with
$b_j^2=|dh_j(p_j)|^2\to\infty$, $b_j\geq1$. Use the conformal
sphere chart centered at $p_j$, normalized so its round factor is
$\lambda(z)=16/(|z|^2+4)^2$, and put
\[
 s_j=b_j^{-1},\qquad
 v_j(z)=h_j\bigl(\sigma_{p_j}(s_jz)\bigr).
\]
The half-energy density $e(v_j)=\tfrac12\sum_i
g(\partial_i v_j,\partial_i v_j)$ satisfies
\[
 e(v_j)(0)=\tfrac12,\qquad
 0\leq\frac{2e(v_j)(z)}{\lambda(s_jz)}\leq1,
 \qquad e(v_j)(z)\leq\tfrac12.
\]
Conformal reparametrization preserves total energy:
$\int_{\mathbb R^2}e(v_j)=E(h_j)$.

The last bound and compactness of $X$ give a locally uniform
subsequence limit. To improve it to first-derivative convergence,
work near a point in a common target chart, available from the
uniform limit. A smooth embedding of $X$ may be chosen so that these
chart coordinates are fixed linear projections of its coordinate
functions near the limiting image. Rescale a small source disk to a
unit disk. The weighted equation, normalized by the target metric,
has a uniformly bounded lower term and a principal error tending to
zero with $\alpha_j-1$. The interior near-Laplacian estimate gives
a common $1/4$-Holder modulus for the first derivatives on the
smaller disk. Arzela--Ascoli for values and first derivatives,
followed by a diagonal selection over compact disks, gives one
target-valued $C^1$ limit $v$ and one subsequence converging in both
values and first derivatives locally uniformly.

It remains to check the equation at zero as well as nonzero
gradients. In the rescaled chart its weight is
\[
 w_j(z)=\left(s_j^2+
 \frac{\sum_i g_{v_j}(\partial_i v_j,\partial_i v_j)}
      {\lambda(s_jz)}\right)^{\alpha_j-1}.
\]
Where $dv\ne0$, the base in this expression has a positive limit,
so $w_j\to1$. Where $dv=0$, that conclusion for the weight alone
need not hold, but the first-variation integrand, linear or
quadratic in $dv_j$, tends to zero. The displayed normalization
bounds $w_j\leq2$; the chart coefficients, derivatives, and each
fixed supported test function give a common integrable bound.
Dominated convergence of the weighted first variation therefore
gives the unweighted harmonic equation everywhere in the weak
sense. The local regularity argument at $\alpha=1$ makes $v$
smooth and harmonic. At the origin its energy density is $1/2$,
so it is nonconstant. On each disk the density integrals converge;
monotone exhaustion gives $E(v)\leq m$.

Conformal inversion makes infinity a finite-energy puncture.
The removable-singularity theorem for a harmonic map from a
punctured surface disk gives a smooth harmonic sphere $u$ with
$E(u)=E(v)$; see
\cite[Theorem 3.6, pp.~13--15]{SacksUhlenbeck1981}.
The extraction retains the subsequence of the original $h_j$,
the centers and positive scales, and convergence of disk areas
from convergence of the actual first derivatives. These facts
are needed to prove that $u$ is non-null. Nonconstancy alone
would not prove it.

Suppose, for a contradiction, that this rescaled $u$ is null-homotopic.
Its area $a$ is positive: harmonicity makes the Hopf differential
holomorphic, hence zero on $S^2$; the resulting nonconstant weakly
conformal sphere has only finitely many branch points. Choose a planar
disk on which the area $D$ of $u$ exceeds $3a/4$; its complementary
tail has area less than $a/4$. For large $j$, the corresponding rescaled
disk of $h_j$ has area greater than $D-a/12$, and
$\operatorname{Area}(h_j)<m+a/12$.

Here are the geometric and topological details of the replacement.
Reflect the complement of the radius-$R$ disk conformally into that
disk. Composing this cap parametrization with $u$ gives a smooth
map $a_R$ whose boundary agrees with $u\circ\sigma$ on $|z|=R$,
and whose disk area is the complementary area of $u$.
In a collar $R-d<|z|<R+d$, interpolate between $a_R$ on the
inside and $v_j$ on the outside by the local contraction of nearby
target pairs. Its spatial derivatives are bounded by a fixed
constant $B$. Its time derivative can be made as small as any
prescribed $A>0$ by requiring the pair to lie sufficiently near
the diagonal, because the contraction fixes diagonal pairs.
First-derivative convergence bounds the derivatives of $v_j$ and
the fixed cap by a common $L$. A smooth radial transition with
derivative at most $H/(2d)$ then gives the collar-area estimate
\[
 4\pi R\left(\frac{HA}{2}BL+d(BL)^2\right).
\]
Choose $A$ first and then $d$ small, and finally $j$ large. This
order makes the collar area less than $a/12$ while preserving a
$C^1$ map across both collar boundaries. Positive rescaling back
to the original sphere preserves these area integrals.

The homotopy class is controlled separately. A null homotopy of
$u$ gives an extension $U$ to the closed unit ball. If $a,b$ are
the two sphere-valued cap parametrizations, then
$U((1-t)a+tb)$ is a homotopy of their images that fixes their
common boundary circle. Compactness of that circle times the
time interval gives a single sufficiently thin collar on which
this entire homotopy remains close to $u\circ\sigma$. First
interpolate $v_j$ to $u\circ\sigma$ on the inside, then perform
the cap homotopy there, using the same target contraction to join
to the unchanged $v_j$ outside. This homotopy is stationary
outside the collar and extends over the sphere's omitted pole.
The resulting replacement sphere is therefore homotopic to the
original $h_j$ and remains non-null. Area bookkeeping subtracts
the original radius-$R$ disk, adds the complementary cap, and
adds the collar error. Its area is bounded above by
\[
 \operatorname{Area}(h_j)-(D-a/12)+a/4+a/12
 <m+\frac a{12}-\frac{3a}{4}+\frac a{12}+\frac a4+\frac a{12}
 =m-\frac a4,
\]
contrary to the definition of $m$. Notice that the homotopy class retained
here is that of $h_j$, not that of a possibly null weak limit.

Consequently $u$ is non-null. Weak conformality gives
$m\leq\operatorname{Area}(u)=E(u)\leq m$, so equality holds.
Small variations of $u$ stay in its free homotopy class, and least energy
therefore implies energy stationarity. The harmonic regularity and
branch-point conclusions give the required smooth branched minimal
sphere, with $m=\operatorname{Area}(u)>0$. Every $C^1$ sphere of area
less than $m$ is null by the definition of this infimum.
\source{MT2007}\dcref{Lemma 18.10 and Claims 18.12--18.13, pp. 424--428}
\source{SacksUhlenbeck1981}\dcref{Theorem 1.6, Proposition 2.3, Main Estimate 3.2, Theorems 3.6, 4.4, 4.6 and 4.7, pp. 5--7, 11--18}

\subsection{Family and class widths}

For a raw null family $F:S^2\to\Lambda^1M$ define
\[
 W_g(F)=\sup_{c\in S^2}A_g(F(c)).
\]
Continuity of $A_g$ on null loops and compactness of $S^2$ imply that this
supremum is finite and attained at some parameter $c_0$. The family width is a
maximum, even though the disk infimum inside $A_g$ need not be a minimum.

The free class range is
\[
 {\cal R}_g(F)=\{W_g(G):G:S^2\to\Lambda^1M\text{ raw null and }G\simeq F\},
\]
where the homotopy is unrestricted in the free loop space, and the free-class
width is $W_g[F]=\inf\mathcal R_g(F)$. The range is nonempty because it
contains $W_g(F)$, and every element is nonnegative. Thus $W_g[F]\geq0$;
it is at most every competitor, and for every $\epsilon>0$ there is
$G\simeq F$ with $W_g(G)<W_g[F]+\epsilon$. If $F\simeq H$,
concatenating homotopies gives equal ranges and $W_g[F]=W_g[H]$.
No minimizing family is asserted.

For $\alpha\in\pi_2(\Lambda^1M,\gamma_x)$, choose a normalized regular
family $\Gamma$ with class $(x,\alpha)$. A raw family $F$ represents
$\alpha$ if it is freely homotopic to $\Gamma$. Define
\[
 W_g(x,\alpha)=\inf\{W_g(F):F\text{ raw null and represents }\alpha\}.
\]
The coherent identification theorem proves that this range equals the free
class range of any representing anchor. Consequently $W_g(x,\alpha)$ is
nonnegative and has the same near-minimizers, independent of the normalized
seed.

\begin{proposition}[Finite nonnegative widths]
\label{prop:filling-width-properties}
\uses{prop:filling-short-area,prop:filling-loop-identification}
Under the compactness, connectedness, smoothness and $\pi_2(M,x)=0$ hypotheses,
the family maximum, free-class infimum and based-class infimum satisfy the
attainment, boundedness, nonnegativity, homotopy invariance and
near-minimizer properties above. Separately, for a compact smooth target
with nontrivial $\pi_2$, the least non-null sphere area is positive and
attained by a branched minimal sphere, as proved above.
\source{MT2007}\dcref{Definition 18.17 and Claim 18.16, p. 430}
\dcref{corrected Corollary 18.28, p. 434}
\lean{PoincareMT.m61Widths}
\end{proposition}

These are three different extremal operations: disk areas are infimized,
parameters of one family are maximized, and representative families are
then infimized. The loop-class identification is used only in the last
operation, to identify the raw free-class range with the based-class range.

\section{Variation of sphere area under Ricci flow}
\label{sec:filling-sphere-variation}

Let $g(t)$ be a smooth Ricci flow on a fixed $n$-manifold for
$a\leq t\leq b$, with $a<b$, and let $f:S^2\to X$ be a fixed $C^1$
map. Suppose $|\operatorname{Ric}_{g(t)}|\leq D$ everywhere on the slab,
where $D\geq0$. Put $A(t)=\operatorname{Area}_{g(t)}(f)$. At a point
where $df$ has rank two, let $G(t)$ be its two-by-two Gram matrix and
$R(t)$ the matrix of the Ricci tensor on the same two derivative vectors.
Since $G'=-2R$, differentiating the square root of the determinant gives
\[
 \frac{d}{dt}\sqrt{\det G}
 =\frac12\operatorname{tr}(G^{-1}G')\sqrt{\det G}
 =-\operatorname{tr}(G^{-1}R)\sqrt{\det G}.
\]
At a rank-deficient point the derivative vectors are linearly dependent
for every metric. The area density is identically zero in time there,
so its derivative and its Ricci-trace density are both zero. Thus the
formula covers the entire parameter plane, including a map of varying rank.

An orthonormal basis of the image two-plane gives
$|\operatorname{tr}(G^{-1}R)|\leq2D$. In particular, the coarser bound
used throughout the construction is
$|\partial_t J(t,z)|\leq4D J(t,z)$. The one-variable exponential
comparison applied pointwise gives
\[
 J(t,z)\leq e^{4D(b-a)}J(a,z),\qquad a\leq t\leq b.
\]
The right side is integrable: the sphere map is $C^1$ on a compact
domain, and its stereographic area density is integrable. It also gives
the integrable derivative bound
$4D e^{4D(b-a)}J(a,z)$ on the entire slab. The mean value inequality
bounds time difference quotients by this same function. Dominated
convergence therefore permits differentiation of the area integral,
including the one-sided derivatives at the slab endpoints. Consequently
\[
 A'(t)=-\int_{\mathbb R^2}\mathcal T(t,z)\,dz,
 \qquad |A'(t)|\leq4D A(t),
\]
where $\mathcal T$ is the inverse-Gram Ricci contraction times area
density, extended by zero at rank-deficient points. In particular,
$e^{-4D(t-a)}A(t)$ is nonincreasing, and for any $s,t\in[a,b]$,
\[
 e^{-4D|t-s|}A(s)\leq A(t)\leq e^{4D|t-s|}A(s).
\]
The variation contains one negative Ricci trace. The factor $1/2$ in
the determinant derivative cancels the factor two in the Ricci-flow
equation; the coarser exponential constant four does not alter this
identity. This corrects the variation line in
\cite[Claim 18.13, p.~427]{MT2007}.

There is a stronger upper bound at a time $t_0$ when $f$ is a branched
minimal sphere in a three-manifold. We give the treatment of branches
explicitly. With respect to the unit round metric $g_0$ write
$f^*g(t_0)=qg_0$, where $q\geq0$ is smooth and positive away from the
finite branch set. Define the smooth functions on the whole sphere
\[
 \tau=\operatorname{tr}_{g_0}f^*\operatorname{Ric}_{g(t_0)},\qquad
 k=\tau-\tfrac12(R_{g(t_0)}\circ f)q.
\]
Both are defined even at branch points, without choosing a normal.
On the regular locus the dimension-three curvature identity identifies
$k$ with $q$ times the ambient sectional curvature of the tangent plane.
Energy stationarity gives the harmonic equation in each target chart.
Its Bochner calculation, together with weak conformality, yields
\[
 \Delta_{g_0}\log q\geq2-2k\quad(q>0),
 \qquad
 \Delta_{g_0}q\geq\frac{|\nabla q|^2}{q}+2q-2qk.
\]
The constant two is twice the Gaussian curvature of the unit round
sphere. In stereographic coordinates this follows by subtracting the
logarithm of the round area density from the logarithm of the map's
area density; continuity extends the regular-locus inequality across
the omitted chart pole whenever $q$ is positive there.

To integrate without taking the logarithm at a branch point, fix
$\epsilon>0$ and set $r_\epsilon=q/(q+\epsilon)$. The scalar chain rule
and the preceding inequality give
\[
 \Delta_{g_0}\log(q+\epsilon)
 \geq 2r_\epsilon(1-k).
\]
For $q>0$, the extra gradient term is
$\epsilon|\nabla q|^2/[q(q+\epsilon)^2]\geq0$.
At a zero of the nonnegative smooth function $q$, its gradient is zero
and its Laplacian is nonnegative, so the same inequality holds there.
The integral of this smooth Laplacian on the closed round sphere is
zero. It follows that
\[
 \int_{S^2}r_\epsilon\,d\mu_0
 \leq\int_{S^2}r_\epsilon k\,d\mu_0.
\]
Now $0\leq r_\epsilon\leq1$ and $r_\epsilon\to1$ away from the
finite branch set, which has measure zero. Dominated convergence, using
integrability of $k$, gives $4\pi\leq\int k\,d\mu_0$.

If $R_{g(t_0)}\geq\rho$, with no restriction on the sign of $\rho$,
then $q\geq0$ and the definitions imply
\[
 \int_{S^2}\tau\,d\mu_0
 =\int_{S^2}\left(k+\tfrac12(R_{g(t_0)}\circ f)q\right)d\mu_0
 \geq4\pi+\frac\rho2\operatorname{Area}_{g(t_0)}(f).
\]
Changing from round coordinates to the stereographic integral identifies
this with the Ricci-trace integral in the fixed-map formula. Compactness
of a flow slab supplies a uniform Ricci bound when needed to apply that
formula. We have proved
\begin{proposition}[Minimal-sphere area variation]
\label{prop:filling-minimal-sphere-variation}
\uses{prop:rf-evolution}
For a fixed smooth sphere map that is branched minimal at $t_0$ in a
compact three-dimensional Ricci flow,
\[
 \left.\frac{d}{dt}\operatorname{Area}_{g(t)}(f)\right|_{t=t_0}
 \leq-4\pi-\frac\rho2\operatorname{Area}_{g(t_0)}(f)
\]
for every scalar lower bound $\rho$ at that time. In particular one may
take the attained minimum of scalar curvature on the compact slice.
At a slab endpoint the derivative is taken within the slab.
\source{MT2007}\dcref{Claims 18.12--18.13, pp. 426--428}
\lean{PoincareMT.m60AreaAndFilling}
\end{proposition}

\section{One initial class for all target components}
\label{sec:filling-fixed-initial-class}

Fix the global flow of a compact simply connected initial three-manifold.
Its initial diffeomorphism transports simple connectedness to $M_0$.
Choose once a smooth component model $C_0$ of $M_0$, with inclusion
$\iota_0:C_0\to M_0$ and basepoint $x_0$.  Since $M_0$ is connected,
$\iota_0(C_0)=M_0$.

For each $T\geq0$ and $x\in M_T$, the ancestry construction gives a path
of components ending in the component containing $x$.  Its initial model
is the same $C_0$, including the selected point $x_0$.  Here is how that
agreement is achieved.  At time zero use $C_0$ in every trace; at a positive
time use the open component model at the trace's selected point.  At zero
the range is the entire connected slice, so this replacement represents
exactly the original trace's component.  A regular-slab diffeomorphism
preserves components and restricts to these models.  At an event, a point
in a surviving outgoing component can be joined within that component to
the retained region, whose inverse retention map supplies an incoming
point.  To find such an interior retained point, a cap's connected local
ball model also contains a negative half-neck, which pulls back into the
retained interior.  If a retained witness lies on the boundary, the
event's boundary correspondence places it on a cap, and the same local
model supplies an interior witness in its component.  Induction through
the finitely many events constructs the trace.
Consequently its component at zero is literally the chosen model, rather
than a model whose basepoint would have to be chosen again for each target.

For a fixed target slice, the terminal component ranges form an open cover.
Compactness gives finitely many of them.  In each selected range choose
the image of its model's basepoint and use the ancestry path assigned to
that terminal point.  The new and old ranges are connected components
containing the same point, so they agree.  This gives a finite cover by
the actual selected paths while preserving their common initial model.
An empty target needs no such paths.

Restrict the time-zero metric by
\[
 g_{C_0}(v,w)=g(0)(d\iota_0(v),d\iota_0(w)).
\]
Corollary~\ref{cor:foundations-homotopy-groups} for a closed simply
connected three-manifold gives
at some point $y\in C_0$ both $\pi_2(C_0,y)=0$ and an isomorphism
$e:\pi_3(C_0,y)\to\mathbb Z$.  Choose a path $p$ in $C_0$ from $y$ to
$x_0$ and put
\[
 \xi_0=e^{-1}(1),\qquad
 \xi=p_*\xi_0,\qquad
 \alpha_0=\Phi_{C_0,x_0}^{-1}(\xi).
\]
The inverse path shows that $p_*$ is injective and preserves zero, so
$\xi\ne0$.  The same inverse-path argument transports the triviality of
$\pi_2$ to $x_0$.  Hence the displayed loop-space identification is
applicable and $\alpha_0$ is nonzero.  The quotient parameter, identification
system and path-transport system used here are fixed together.

Now define the single real number
\begin{equation}
 W_0=W_{g_{C_0}}(x_0,\alpha_0)\geq0.
 \label{eq:filling-fixed-initial-width}
\end{equation}
It is finite by Proposition~\ref{prop:filling-width-properties}.  All
choices in this formula precede the choice of $T$ or of a target point.
For each target path, use the equality of its zero-time component with
$C_0$ to supply this same metric and class to the width-transport theorem.
That theorem retains the initial metric and class, so the width function
along every such path satisfies
\[
 W_{T,x}(0)=W_0.
\]
Thus the extinction comparison can choose its observation time from one
finite initial width and then apply to every component reaching that time.
No positive lower bound for $W_0$ is asserted or needed at this interface.
The construction is the simply connected, time-zero specialization of
\cite[Proposition 18.18 and the proof of Theorem 18.1,
pp.~431--432]{MT2007}.

\section{Interfaces to the extinction argument}
\label{sec:filling-width-interfaces}

The next chapters receive a fixed coherent label on one actual surgery
component, a finite nonnegative width, and arbitrarily close raw
representatives. The ramp chapter replaces a raw family by a polygonal family
while preserving its free class; the annular-comparison chapter deforms that
family with a controlled area error. The extinction chapter transports the
same class across regular intervals and surgery events, using path-based
basepoint maps. The scalar-clock argument compares this finite width with the
negative extinction profile.

The following distinctions remain visible: sphere separation concerns the
actual complement; factor injection proves triviality only after its parent
group and survivor guard are supplied; the short-loop theorem gives an actual
disk while filling area is an infimum; $W_g(F)$ is a compact-parameter maximum
whereas $W_g[F]$ and $W_g(x,\alpha)$ are infima; and basepoint changes are induced
by explicit paths and whiskering, not arbitrary group isomorphisms.
